% !TEX root = ../main.tex
%%
% BIThesis 本科毕业设计论文模板 —— 使用 XeLaTeX 编译
%%

% 中英文摘要章节
\begin{abstract}
随着软件系统规模与复杂度持续增长，软件质量保障面临严峻挑战。传统软件测试高度依赖人工经验，存在效率低、覆盖不足、成本高等问题。如何提高测试效率、保证测试质量，已成为软件工程领域的重要研究问题。

近年来，敏捷开发与持续集成模式得到广泛应用，代码变更更加频繁，回归测试成为保证软件质量的关键环节。传统的回归测试主要依赖人工编写测试用例，难以快速响应频繁的代码变更。

代码变更影响分析是回归测试的核心环节，其目标在于识别代码变更可能影响的系统范围，从而指导测试资源的分配。传统方法仅能识别行级变化，难以反映代码的真实语义与影响范围。

大语言模型在代码理解、自然语言处理等方面展现出强大能力，多智能体系统通过多个Agent的协同决策，能够处理复杂多步骤的任务。本研究将大语言模型与多智能体系统相结合，提出基于多智能体的软件测试自动化方法，实现从代码变更到测试用例生成的端到端自动化。

本研究的主要贡献包括：（1）提出基于结构化中间表示的代码变更影响分析方法；（2）设计多智能体协同测试生成框架；（3）构建规则约束与LLM生成相结合的混合测试策略；（4）实现基于TestRepairAgent的测试修复机制。

\end{abstract}

% 英文摘要章节
\begin{abstractEn}
With the continuous growth of software system scale and complexity, software quality assurance faces severe challenges. Traditional software testing heavily relies on manual experience, resulting in low efficiency, insufficient coverage, and high costs.

In recent years, agile development and continuous integration have been widely adopted, making code changes more frequent and regression testing critical for software quality. Traditional regression testing mainly depends on manually编写测试用例，making it difficult to respond quickly to frequent code changes.

Change impact analysis is a core component of regression testing, aiming to identify the system scope affected by code changes. Traditional methods can only identify line-level changes and fail to reflect the true semantics and impact scope.

Large language models demonstrate strong capabilities in code understanding and natural language processing. Multi-agent systems can handle complex multi-step tasks through collaborative decision-making. This research combines large language models with multi-agent systems to propose an automated software testing method, achieving end-to-end automation from code changes to test case generation.

\end{abstractEn}
